\section{Experimentopzet}
\label{sec:experimentopzet}

Om de werking van RustProv te evalueren, werd een testomgeving opgezet die aansluit bij een realistische onderwijscontext. Daarbij werd gekozen voor twee afzonderlijke scenario's: een Windows-omgeving en een Linux-omgeving. Beide omgevingen zijn gebaseerd op opstellingen uit het vak SDP, zodat de experimenten zoveel mogelijk aansluiten bij een bestaande en relevante praktijksituatie binnen de opleiding.

De keuze voor deze opzet maakt het mogelijk om RustProv te testen in scenario's met meerdere virtuele machines, onderlinge afhankelijkheden en afzonderlijke provisioningscripts per machine. Op die manier kan niet alleen worden nagegaan of de tool technisch correct werkt, maar ook of de gekozen aanpak bruikbaar is in een onderwijs- en labo-omgeving.

\subsection{Windows-omgeving}

Voor het Windows-scenario werd een omgeving uitgewerkt die gebaseerd is op een SDP-opstelling met meerdere samenwerkende systemen binnen één netwerkstructuur. Deze omgeving werd gebruikt om na te gaan of RustProv ook in een meer complexe Windows-context correct kan functioneren.

\subsubsection{Enkele virtuele machine}

Als eerste testgeval werd één afzonderlijke Windows-VM gebruikt. Deze test had als doel na te gaan of RustProv in staat is een enkele Windows-machine correct aan te maken, op te starten en te provisionen. In deze opstelling werd een Windows-client gebruikt waarop de nodige beheertools, waaronder RSAT, via een provisioningscript werden voorbereid. Deze configuratie vormt een beperkt maar representatief scenario voor het evalueren van de basisfunctionaliteit van de provisioner binnen een Windows-omgeving.

\paragraph{Script 1}
Het derde 
\begin{lstlisting}[language=PowerShell,caption={Provisioningscript voor de Windows-client.}]
    % script 3 hier invoegen
\end{lstlisting}

\subsubsection{Meerdere virtuele machines}

Naast het enkelvoudige testgeval werd ook een uitgebreidere Windows-opstelling gebruikt met drie virtuele machines: twee servers en één client. Deze configuratie is gebaseerd op een SDP-scenario waarin meerdere Windows-systemen onderling samenwerken. Daardoor kan niet alleen de provisioning van afzonderlijke machines worden getest, maar ook de opstartvolgorde en de samenhang tussen verschillende systemen binnen één omgeving.

Voor elke virtuele machine werd een afzonderlijk provisioningscript voorzien. Op die manier kon elke machine een eigen rol en configuratie krijgen, terwijl de provisioner toch één samenhangend scenario moest uitvoeren. Deze opdeling maakt het mogelijk om per systeem te controleren of de toegewezen configuratie correct werd toegepast.

\paragraph{Script 1}
Het eerste 

\begin{lstlisting}[language=PowerShell,caption={Provisioningscript voor Windows Server 1.}]
    % script 1 hier invoegen
\end{lstlisting}

\paragraph{Script 2}
Het tweede script 

\begin{lstlisting}[language=PowerShell,caption={Provisioningscript voor Windows Server 2.}]
    % script 2 hier invoegen
\end{lstlisting}

\paragraph{Script 3}
Het derde 
\begin{lstlisting}[language=PowerShell,caption={Provisioningscript voor de Windows-client.}]
    % script 3 hier invoegen
\end{lstlisting}

\subsection{Linux-omgeving}

Voor het Linux-scenario werd een opstelling voorzien met drie virtuele machines op basis van AlmaLinux. Deze omgeving is gebaseerd op een configuratie uit het vak SDP en werd gekozen omdat ze representatief is voor een typische Linux-opstelling binnen de opleiding. Daardoor sluit dit experiment aan bij de concrete onderwijscontext waarvoor RustProv bedoeld is.

\subsubsection{Enkele virtuele machine}

Als eerste Linux-testgeval werd één afzonderlijke AlmaLinux-VM gebruikt. Deze test had als doel te verifiëren of RustProv een enkele Linux-machine correct kan opstarten en provisionen. In deze configuratie werd een eenvoudige serveropstelling gebruikt waarbij onder andere Apache werd voorbereid via een provisioningscript. Dit beperkte scenario laat toe de basiswerking van de Linux-provisioning afzonderlijk te evalueren. 


\paragraph{Script 1}
Het eerste script 
\begin{lstlisting}[language=bash,caption={Provisioningscript voor Linux VM 1.}]
    # script 1 hier invoegen
\end{lstlisting}

\paragraph{Script 2}
Het tweede script 

\subsubsection{Meerdere virtuele machines}

Daarnaast werd ook een Linux-opstelling met drie virtuele machines getest. Deze configuratie sluit aan bij een SDP-omgeving waarin meerdere AlmaLinux-systemen gelijktijdig worden ingezet. Het gebruik van meerdere Linux-VM's maakt het mogelijk om RustProv te evalueren in een scenario met meerdere nodes, afzonderlijke scripts en een gezamenlijke uitvoeringscontext.

Voor elke Linux-VM werd een apart provisioningscript voorzien. Daardoor kon elke machine een eigen configuratie en rol krijgen binnen de bredere testomgeving. Deze aanpak maakt het mogelijk om te beoordelen of RustProv de juiste scriptstappen op de juiste machine uitvoert en of meerdere Linux-systemen correct na elkaar of gezamenlijk verwerkt kunnen worden. 

\subsubsection{Linux-script 1: database en Loki}

Dit script installeert de databaseomgeving en Loki, zodat loggegevens centraal kunnen worden opgeslagen en later geanalyseerd. Het vormt de basislaag van de monitoringopstelling.

\begin{lstlisting}[language=bash,caption={Linux-script 1: installatie van database en Loki.}]
    #!/bin/bash
    
    # Install dependencies
    sudo yum update -y
    sudo yum install -y wget curl unzip tar
    
    # Download and install Loki
    cd /opt
    sudo mkdir -p /opt/loki
    cd /opt/loki
    sudo wget https://github.com/grafana/loki/releases/download/v2.9.2/loki-linux-amd64.zip
    sudo unzip loki-linux-amd64.zip
    sudo chmod +x loki-linux-amd64
    sudo mv loki-linux-amd64 /usr/local/bin/loki
    
    # Create Loki config
    sudo mkdir -p /etc/loki
    cat <<EOF | sudo tee /etc/loki/loki-config.yaml > /dev/null
    auth_enabled: false
    
    server:
    http_listen_port: 3100
    
    common:
    path_prefix: /loki
    storage:
    filesystem:
    chunks_directory: /loki/chunks
    rules_directory: /loki/rules
    replication_factor: 1
    ring:
    kvstore:
    store: inmemory
    
    schema_config:
    configs:
    - from: 2024-01-01
    store: tsdb
    object_store: filesystem
    schema: v13
    index:
    prefix: loki_index_
    period: 24h
    
    storage_config:
    tsdb_shipper:
    active_index_directory: /loki/index
    cache_location: /loki/index_cache
    EOF
    
    # Create directories
    sudo mkdir -p /loki/chunks /loki/rules /loki/index /loki/index_cache
    
    # Create systemd service
    cat <<EOF | sudo tee /etc/systemd/system/loki.service > /dev/null
    [Unit]
    Description=Loki
    After=network.target
    
    [Service]
    ExecStart=/usr/local/bin/loki -config.file=/etc/loki/loki-config.yaml
    Restart=always
    User=root
    
    [Install]
    WantedBy=multi-user.target
    EOF
    
    sudo systemctl daemon-reload
    sudo systemctl enable --now loki
\end{lstlisting}

\subsubsection{Linux-script 2: WordPress en Promtail}

Dit script installeert WordPress en Promtail. WordPress levert de webtoepassing, terwijl Promtail de logs van het systeem doorstuurt naar Loki voor centrale logverzameling.

\begin{lstlisting}[language=bash,caption={Linux-script 2: installatie van WordPress en Promtail.}]
    #!/bin/bash
    
    # Update system
    sudo yum update -y
    sudo yum install -y httpd wget php php-mysqlnd php-fpm unzip curl
    
    # Install WordPress
    cd /tmp
    wget https://wordpress.org/latest.tar.gz
    tar -xzf latest.tar.gz
    sudo cp -r wordpress/* /var/www/html/
    
    # Set permissions
    sudo chown -R apache:apache /var/www/html
    sudo chmod -R 755 /var/www/html
    
    # Configure Apache
    sudo systemctl enable httpd
    sudo systemctl start httpd
    
    # Install Promtail
    cd /opt
    sudo mkdir -p /opt/promtail
    cd /opt/promtail
    sudo wget https://github.com/grafana/loki/releases/download/v2.9.2/promtail-linux-amd64.zip
    sudo unzip promtail-linux-amd64.zip
    sudo chmod +x promtail-linux-amd64
    sudo mv promtail-linux-amd64 /usr/local/bin/promtail
    
    # Create Promtail config
    sudo mkdir -p /etc/promtail
    cat <<EOF | sudo tee /etc/promtail/promtail-config.yaml > /dev/null
    server:
    http_listen_port: 9080
    
    positions:
    filename: /tmp/positions.yaml
    
    clients:
    - url: http://127.0.0.1:3100/loki/api/v1/push
    
    scrape_configs:
    - job_name: system
    static_configs:
    - targets:
    - localhost
    labels:
    job: varlogs
    __path__: /var/log/*log
    EOF
    
    # Create systemd service
    cat <<EOF | sudo tee /etc/systemd/system/promtail.service > /dev/null
    [Unit]
    Description=Promtail
    After=network.target
    
    [Service]
    ExecStart=/usr/local/bin/promtail -config.file=/etc/promtail/promtail-config.yaml
    Restart=always
    User=root
    
    [Install]
    WantedBy=multi-user.target
    EOF
    
    sudo systemctl daemon-reload
    sudo systemctl enable --now promtail
\end{lstlisting}

\subsubsection{Linux-script 3: monitoring met Grafana}

Dit script installeert Grafana en maakt de monitoringomgeving volledig. Grafana gebruikt de data uit Loki om dashboards en logoverzichten weer te geven.

\begin{lstlisting}[language=bash,caption={Linux-script 3: installatie van Grafana en monitoring.}]
    #!/bin/bash
    
    # Install Grafana repository
    sudo yum install -y https://dl.grafana.com/enterprise/release/grafana-enterprise-11.1.0-1.x86_64.rpm
    
    # Enable and start Grafana
    sudo systemctl daemon-reload
    sudo systemctl enable --now grafana-server
    
    # Add Loki datasource via provisioning
    sudo mkdir -p /etc/grafana/provisioning/datasources
    cat <<EOF | sudo tee /etc/grafana/provisioning/datasources/loki.yaml > /dev/null
    apiVersion: 1
    
    datasources:
    - name: Loki
    type: loki
    access: proxy
    url: http://127.0.0.1:3100
    isDefault: true
    EOF
    
    sudo systemctl restart grafana-server
\end{lstlisting}